\subsection{Universality of the Saturation Scale}
  \label{subsec:universality}

  A key result of the analysis is the stability of the saturation scale $a_\star$ across the sample.
  Allowing moderate variations of $a_\star$ does not significantly improve the
  quality of the fits and tends to introduce degeneracies with baryonic mass normalization.

  This supports the interpretation of $a_\star$ as a universal effective scale
  rather than a phenomenological fitting parameter.
  All galaxy-to-galaxy variability is accounted for by the observed baryonic structure.

  The absence of halo-by-halo tuning distinguishes this approach from standard
  dark matter halo modeling, which typically requires galaxy-dependent halo profiles and parameters.

  The universality of the saturation scale follows from the existence of a
  fundamental bound on the effective flux that can be supported by the geometric substrate.
  Unlike dark matter halo models, the present framework involves a single
  acceleration scale set by an intrinsic saturation of the underlying geometric response.
